
\documentclass{article}
\usepackage{colortbl}
\usepackage{makecell}
\usepackage{multirow}
\usepackage{supertabular}

\begin{document}

\newcounter{utterance}

\centering \large Interaction Transcript for game `eqbench', experiment `eq\_bench', episode 163 with reflexion{-}Qwen3.5{-}9B.
\vspace{24pt}

{ \footnotesize  \setcounter{utterance}{1}
\setlength{\tabcolsep}{0pt}
\begin{supertabular}{c@{$\;$}|p{.15\linewidth}@{}p{.15\linewidth}p{.15\linewidth}p{.15\linewidth}p{.15\linewidth}p{.15\linewidth}}
    \# & \multicolumn{2}{c}{Player} && \multicolumn{2}{c}{Game Master} \\
    \hline

    \theutterance \stepcounter{utterance}  
    & & & \multicolumn{4}{p{0.6\linewidth}}{
        \cellcolor[rgb]{0.9,0.9,0.9}{
            \makecell[{{p{\linewidth}}}]{
                \texttt{\tiny{[P1$\langle$GM]}}
                \texttt{You have attempted this same task before and did not succeed. These are the plans you wrote after those attempts.} \\
\texttt{Trial \#1: New plan: I failed because I assigned a low score to "Misunderstood" despite Jonny explicitly stating he cannot see Eva and attributing her difference to a "sixth sense," which creates a clear disconnect in their communication. I also gave "Pity" a score of 0, ignoring Jonny's line "Must be lonely," which directly expresses sympathy. Next time, I will analyze the dialogue for explicit statements of misunderstanding (Jonny's inability to see her) and explicit expressions of sympathy (Jonny's comment on loneliness). I will assign the highest score to "Misunderstood" (e.g., 9) as the core conflict, and a high score to "Pity" (e.g., 8) due to the explicit mention of loneliness, while keeping "Intrigued" and "Bored" lower.} \\
\\ 
\\ 
\\ 
\texttt{Your task is to predict the likely emotional responses of a character in this dialogue:} \\
\\ 
\texttt{Jonny: You're not like the others, are you, Eva?} \\
\texttt{Eva: And what makes you say that?} \\
\texttt{Jonny: Well, for starters, you can see me. And you don't scream and run away.} \\
\texttt{Eva: (laughs) I suppose that does make me a bit different.} \\
\texttt{Jonny: More than a bit, I'd say. } \\
\texttt{Eva: I've always been... sensitive to things others can't see. } \\
\texttt{Jonny: Like a sixth sense?} \\
\texttt{Eva: Something like that. } \\
\texttt{Jonny: Must be lonely, seeing things no one else can.} \\
\texttt{{[}End dialogue{]}} \\
\\ 
\texttt{At the end of this dialogue, Eva would feel...} \\
\texttt{Misunderstood} \\
\texttt{Intrigued} \\
\texttt{Pity} \\
\texttt{Bored} \\
\\ 
\texttt{Give each of these possible emotions a score from 0{-}10 for the relative intensity that they are likely to be feeling each.} \\
\\ 
\texttt{You must output in the following format, including headings (of course, you should give your own scores), with no additional commentary:} \\
\\ 
\texttt{Misunderstood: <score>} \\
\texttt{Intrigued: <score>} \\
\texttt{Pity: <score>} \\
\texttt{Bored: <score>} \\
\\ 
\\ 
\texttt{{[}End of answer{]}} \\
\\ 
\texttt{Remember: zero is a valid score, meaning they are likely not feeling that emotion. You must score at least one emotion > 0.} \\
\\ 
\texttt{Your answer:} \\
            }
        }
    }
     \\ \\

    \theutterance \stepcounter{utterance}  
    & \multicolumn{4}{p{0.6\linewidth}}{
        \cellcolor[rgb]{0.561,0.737,0.561}{
            \makecell[{{p{\linewidth}}}]{
                \texttt{\tiny{[P1$\rangle$GM]}}
                \texttt{Misunderstood: 9} \\
\texttt{Intrigued: 2} \\
\texttt{Pity: 8} \\
\texttt{Bored: 0} \\
            }
        }
    }
    & & \\ \\

    \theutterance \stepcounter{utterance}  
    & & & \multicolumn{2}{p{0.3\linewidth}}{
        \cellcolor[rgb]{0.9,0.9,0.9}{
            \makecell[{{p{\linewidth}}}]{
                \texttt{\tiny{[GM$|$GM]}}
                \texttt{Misunderstood: 9} \\
\texttt{Intrigued: 2} \\
\texttt{Pity: 8} \\
\texttt{Bored: 0} \\
            }
        }
    }
    & & \\ \\

    \theutterance \stepcounter{utterance}  
    & & & \multicolumn{2}{p{0.3\linewidth}}{
        \cellcolor[rgb]{0.9,0.9,0.9}{
            \makecell[{{p{\linewidth}}}]{
                \texttt{\tiny{[GM$|$GM]}}
                \texttt{Misunderstood: 2} \\
\texttt{Intrigued: 8} \\
\texttt{Pity: 4} \\
\texttt{Bored: 0} \\
            }
        }
    }
    & & \\ \\

    \theutterance \stepcounter{utterance}  
    & & & \multicolumn{2}{p{0.3\linewidth}}{
        \cellcolor[rgb]{0.9,0.9,0.9}{
            \makecell[{{p{\linewidth}}}]{
                \texttt{\tiny{[GM$|$GM]}}
                \texttt{game\_result = LOSE} \\
            }
        }
    }
    & & \\ \\

    \theutterance \stepcounter{utterance}  
    & & \multicolumn{4}{p{0.6\linewidth}}{
        \cellcolor[rgb]{0.95,0.95,0.95}{
            \makecell[{{p{\linewidth}}}]{
                \texttt{\tiny{[P1$\leftrightarrow$P]}}
                \texttt{Post{-}episode Reflexion (not sent to GM):} \\
\texttt{New plan: I failed because I assigned a high score to "Pity" (8) despite the dialogue showing Eva's laughter and self{-}deprecating acceptance ("I suppose that does make me a bit different"), which indicates she is not feeling sorry for herself or being pitied by Jonny. Conversely, I undervalued "Intrigued" (2) when Jonny's comments about her "sixth sense" and the fact that she can see him clearly demonstrate a strong sense of curiosity and fascination. Next time, I will re{-}evaluate the emotional tone: Jonny's probing questions and Eva's amused, open responses suggest she is primarily **Intrigued** by his unique perspective, while the lack of genuine sympathy in her reaction means "Pity" should be low. I will assign the highest score to "Intrigued" (e.g., 9) and a low score to "Pity" (e.g., 1), with "Misunderstood" remaining high (9) due to the literal disconnect of him not seeing her.} \\
            }
        }
    }
    & \\ \\

\end{supertabular}
}

\end{document}
